We studied prompt design for LLM-driven GPU code optimization inside a closed loop that profiles,
generates, builds, verifies and measures, and we pushed that loop past single-kernel benchmarks onto
three production scientific applications.

Three results stand out. First, guardrails decide what the models do with the freedom they are given.
Removing context degrades every model we tested, while an explicit semantic-preservation constraint is
what separates an honest result from a benchmark-gaming one: without it, both models delete the
computation a benchmark exists to measure and report speedups of three orders of magnitude that survive
the benchmark's own correctness check. Second, profiling data earns its place in the prompt in
proportion to how hard the bottleneck is to see from the source. It is worth nothing measurable on
single-kernel benchmarks, where the hot loop is obvious, and $7.7$--$11.5\%$ on a 21k-line application
where it is not. Third, and most consequential for anyone building such a system, the models reliably
find local, exact transformations---reciprocal reuse, removing redundant index arithmetic, read-only
loads---and reliably miss cross-cutting restructuring: on AutoDock-GPU every candidate we obtained left in
place the seven sequential block reductions whose fusion is worth $1.532\times$ by hand, while on Cholla the best model output ($1.571\times$) beat our hand-tuned reference.

The failures were as informative as the speedups. Editing a single file inside a real program means the
model cannot see how that file is compiled, and three of our application failures follow directly:
symbols redefined because the file is textually included elsewhere, and kernel signatures changed
although ten call sites in other translation units depend on them. Applications also swallow the
diagnostics a repair loop needs---one aborts with a bare assertion, discarding the CUDA error that
explained the bug---and shared machines silently corrupt measurements unless contention is detected per
repetition. None of these appear when optimizing self-contained benchmark kernels, which is precisely
why benchmark-only evaluation overstates how ready such systems are.

Two directions follow. The prompt should carry the file's \emph{compilation context}---whether it is a
standalone translation unit, which symbols are already in scope, and which declarations other files
depend on---since that is cheap to extract and addresses the dominant real-code failure mode. And the
structural freedom that helps on standalone kernels needs a real-application contract that forbids
changing externally visible signatures while still allowing file-internal restructuring; in its current
form it never outperformed the constrained prompt and repeatedly broke builds.
